\chapter{Conclusion}

\addcontentsline{toc}{chapter}{Conclusion}

\textbf{\nomProjet} vise à répondre à un besoin potentiel de la communauté étudiante : faciliter l'achat et la vente de produits de seconde main dans un espace centralisé, accessible et adapté à leurs attentes.

Le projet propose un périmètre réaliste pour cette première version (MVP), comprenant la gestion des utilisateurs, la publication et la consultation des annonces, la recherche et le filtrage, la prise de contact entre acheteurs et vendeurs, ainsi qu'un mécanisme minimal de signalement et de modération. Les paiements en ligne, la livraison intégrée, l'échange sous forme de troc et les fonctionnalités avancées sont, quant à eux, reportés à des versions futures.

La réussite du projet dépendra principalement de la validation des besoins auprès des étudiants, d'une bonne priorisation des fonctionnalités et d'une planification rigoureuse sur les six séances de TP. En se concentrant sur l'essentiel et en restant vigilant vis-à-vis du chemin critique, l'équipe sera en mesure de livrer une première version fonctionnelle et cohérente avec les objectifs fixés.
